\claimheading{Probing Classifiers as Mechanistic Evidence}

\textbf{Source.} \emph{Probing Classifiers: Promises, Shortcomings, and Advances} \citep{belinkov2022probing}.

\textbf{Description.} A classifier $g$ is trained on a model's intermediate output $f_l(x)$ to predict a property $z$, and its accuracy is reported as evidence that the model has learned information relevant to $z$. The audited object is the method rather than one experiment: the anchor is a review that runs no experiments of its own, so the evidence sits in the studies it formalizes. Evaluated on: every architecture family the method has been applied to, from static word embeddings to transformers. Description mode: representational.

\vspace{4pt}
{\footnotesize
\renewcommand{\arraystretch}{1.0}
\begin{longtable}{@{}llll>{\raggedright\arraybackslash}p{4.6cm}@{}}
\caption{\textbf{Probing classifiers.} Full 36-criterion audit. Status: \textbf{C} = Confirmed,
\textbf{PC} = Partially confirmed, \textbf{U} = Untested, \textbf{I} = Inconclusive,
\textbf{D} = Disconfirmed, \textbf{N/A} = Not applicable.}\\
\toprule
\textbf{ID} & \textbf{Criterion} & \textbf{Status} & \textbf{Evidence} \\
\endfirsthead
\toprule
\textbf{ID} & \textbf{Criterion} & \textbf{Status} & \textbf{Evidence} \\
\endhead
\midrule
\multicolumn{4}{@{}l}{\textit{Construct Validity}} \\
C1  & Falsifiability            & C   & The prediction can fail and did, in work the anchor reports side by side \\
C2  & Structural plausibility   & PC  & Every symbol a probing result depends on is named and typed \\
C3  & Convergent validity       & PC  & Probe families agree under accuracy and reverse under other measures \\
C4  & Discriminant validity     & D   & Two neighbors the instrument must separate: probe memorisation, random features \\
C5  & Nomological validity      & PC  & Tied to mutual information, description length and information gain \\
C6  & Complementation validity  & N/A & The audited object is a method, not a construct with named parts \\
\midrule
\multicolumn{4}{@{}l}{\textit{Measurement Validity}} \\
M1  & Reliability               & U   & A review reports no variance; the anchor states what it requires of studies \\
M2  & Baseline separation       & C   & The criterion the literature solved; three control families formalized \\
M3  & Stability                 & I   & Ordering survives under accuracy and inverts under an alternative measure \\
M4  & Calibration               & D   & The shortcomings section opens by denying the measurement has a scale \\
M5  & Sensitivity               & PC  & Skylines and floors exist; no planted property is ever recovered \\
M6  & Invariance                & PC  & Invariance is what probing measures, so the instrument is the question \\
M7  & Selection correction      & U   & Three control families catalogued; the omission from the catalogue is the gap \\
\midrule
\multicolumn{4}{@{}l}{\textit{Internal Validity}} \\
I1  & Necessity                 & I   & Standard probing performs no removal; removal studies disagree four ways \\
I2  & Sufficiency               & N/A & Probing is a map out of the representation, with no path writing back in \\
I3  & Minimality                & U   & The unit is a whole intermediate output; no operation removes part of one \\
I4  & Specificity               & D   & A control dataset holds the property constant and probing fails on it \\
I5  & Rival mechanism exclusion & PC  & Control tasks are solvable only by memorisation, and separate cleanly \\
I6  & Double dissociation       & U   & Needs two properties and two behaviors; the framework supplies one of each \\
I7  & Confound control          & PC  & Probe capacity, information in a random baseline and task difficulty \\
I8  & Confounding sensitivity   & U   & The anchor answers confounding with designed controls, never with a bound \\
I9  & Epistatic interaction     & U   & Non-additivity needs two manipulable components; the object is one output \\
I10 & Rescue reversibility      & U   & Every catalogued intervention runs one way, projecting or training out \\
I11 & Onset coupling            & U   & The original model is taken as given and trained once \\
I12 & Offset coupling           & U   & The converse needs the training sequence onset coupling also lacks \\
\midrule
\multicolumn{4}{@{}l}{\textit{External Validity}} \\
E1  & Intervention reach        & I   & Probing has no intervention of its own; those applied to it disagree \\
E2  & Prompt generalization     & PC  & One line of work reports the same property across several datasets \\
E3  & Cross-task generalization & PC  & Transfers to any property with an annotated dataset, which is the appeal \\
E4  & Cross-model generalization& C   & Portability is established without qualification, from static embeddings on \\
E5  & Graded response           & N/A & No intervention, so no strength to dial and no response to grade \\
E6  & Novel prediction          & PC  & Four design changes followed probing results and then worked \\
\midrule
\multicolumn{4}{@{}l}{\textit{Interpretive Validity}} \\
V1  & Level declaration         & PC  & The formalism fixes which objects a probing result is about \\
V2  & Level-evidence match      & D   & The evidence supports extractability; the claim made is representation \\
V3  & Alternative level         & C   & Every rival reading of a good score is given a measure of its own \\
V4  & Unlicensed labeling       & PC  & ``The model knows z'' projects an epistemic relation onto decodability \\
V5  & Scope declaration         & C   & A scope declaration end to end, conceding limits before any content \\
\midrule
\multicolumn{4}{@{}l}{\textbf{Total:} 5 Confirmed, 12 Partially confirmed, 3 Inconclusive, 9 Untested, 4 Disconfirmed, 3 Not applicable} \\
\multicolumn{4}{@{}l}{\textbf{Verdict: Proposed}} \\
\bottomrule
\end{longtable}}

\clearpage
\begin{table}[H]
\centering
\footnotesize
\caption{\textbf{Probing classifiers.} Readings of the claim, from the version the authors state to the version the
field cites.}
\label{tab:probing-readings}
\renewcommand{\arraystretch}{1.18}
\begin{tabular}{@{}>{\raggedright\arraybackslash}p{3.3cm}l>{\raggedright\arraybackslash}p{4.4cm}@{}}
\toprule
\textbf{Reading} & \textbf{Verdict} & \textbf{Missing} \\
\midrule
\textbf{Primary.} A probing result licenses that property $z$ is extractable from the representation by a classifier of the stated capacity
  & Proposed
  & Necessity: the method performs no removal of its own, and where removal was performed the studies the anchor reports disagree four ways (I1, E1) \\
\addlinespace[2pt]
A good probe score shows that the model represents $z$
  & Disconfirmed
  & Nothing --- tested and failed: control tasks attribute most of a nonlinear probe's accuracy to the probe itself, and even random features decode the property (C4) \\
\addlinespace[2pt]
A good probe score shows that the model uses $z$ in producing its output
  & Disconfirmed
  & Nothing --- tested and failed: a control dataset holds $z$ non-discriminative for the original task and the probe recovers it anyway, so decodability does not localize to the task (I4) \\
\addlinespace[2pt]
Interventions on a probe-identified direction show which features the model uses
  & Underdetermined
  & A result that decides the four-way disagreement among the intervention studies of \S4.3 (I1, E1). This is the broad scope, under which \textbf{I2} and \textbf{E5} are unattempted; under probing as \S2 defines it, a read-out with no path writing back in, they are not applicable \\
\bottomrule
\end{tabular}
\end{table}

\begin{table}[H]
\centering
\footnotesize
\caption{\textbf{Probing classifiers.} Verdict: \textbf{Proposed}. A dash means the tier is reached; a criterion at
Inconclusive, Untested or Disconfirmed blocks promotion.}
\label{tab:probing-verdict}
\renewcommand{\arraystretch}{1.18}
\begin{tabular}{@{}l>{\raggedright\arraybackslash}p{3.6cm}>{\raggedright\arraybackslash}p{4.4cm}@{}}
\toprule
\textbf{Tier} & \textbf{Requires} & \textbf{Missing} \\
\midrule
Proposed & C1--C2; one admissible measurement & --- \\
\addlinespace[2pt]
Causally Suggestive & Necessity (I1); baselines (M2) & \textbf{I1} inconclusive: the anchor sets Giulianelli's positive conclusion against Elazar's negative one and prefers neither \\
\addlinespace[2pt]
Mechanistically Supported & Sufficiency (I2); reach (E1); specificity (I4) & \textbf{I4} Disconfirmed; \textbf{E1} inconclusive; \textbf{I2} not applicable: a read-out has no write path \\
\addlinespace[2pt]
Triangulated & Convergence (C3); replication (E2/E4); dissociation (I6) & \textbf{I6} unattempted: one property and one behavior per experiment; \textbf{C3}: probe families agree under accuracy and reverse under selectivity \\
\addlinespace[2pt]
Validated & Calibration (M1--M7); interpretive audit (V1--V5); reach (E2, E4) & \textbf{M4}, \textbf{V2} Disconfirmed; \textbf{M3} inconclusive; \textbf{M1}, \textbf{M7} unattempted \\
\bottomrule
\end{tabular}
\end{table}
